% Môn: Vật lý
% Lớp: 10
% Chương: Chương I. Mở đầu
% Bài: L10C1 Bài 2. Các quy tắc an toàn trong phòng thực hành Vật lí · An toàn khi sử dụng dụng cụ thủy tinh, vật sắc nhọn
% Số câu: 185
% App đọc file này trực tiếp — sửa rồi Commit trên GitHub, bấm Đồng bộ GitHub .tex

% L10C1 Bài 2. Các quy tắc an toàn trong phòng thực hành Vật lí

% ===== Câu 1 =====
% ID: L10C1B2-01-DS
% Mức: NB
\dangbt{An toàn khi sử dụng dụng cụ thủy tinh, vật sắc nhọn}
\begin{ex}
% ID: L10C1B2-01-DS
% Mức: NB
Trong tình huống sau khi hoàn thành phép đo đối với dụng cụ thủy tinh và vật sắc nhọn, hãy đánh giá các phát biểu sau.
\choiceTF
{\True Cần kiểm tra vết nứt, cầm đúng vị trí và thu gom mảnh vỡ bằng dụng cụ chuyên dụng.}
{Có thể dùng tay không nhặt mảnh kính vỡ nếu thao tác diễn ra nhanh.}
{\True Phải báo cho giáo viên hoặc người phụ trách khi phát hiện nguy cơ vượt khả năng kiểm soát.}
{\True Chỉ tiếp tục thí nghiệm sau khi nguyên nhân nguy hiểm đã được loại bỏ hoặc kiểm soát.}
\loigiai{
\begin{itemchoice}
\item (Đúng) Cần kiểm tra vết nứt, cầm đúng vị trí và thu gom mảnh vỡ bằng dụng cụ chuyên dụng. (Vì): Việc kiểm tra vết nứt, cầm đúng vị trí và thu gom mảnh vỡ bằng dụng cụ chuyên dụng là các biện pháp an toàn cần thiết khi xử lý dụng cụ thủy tinh bị vỡ hoặc vật sắc nhọn
\item (Sai) Có thể dùng tay không nhặt mảnh kính vỡ nếu thao tác diễn ra nhanh. (Vì): Tuyệt đối không được dùng tay không nhặt mảnh kính vỡ, dù thao tác có diễn ra nhanh hay chậm, vì điều này tiềm ẩn nguy cơ gây thương tích nghiêm trọng
\item (Đúng) Phải báo cho giáo viên hoặc người phụ trách khi phát hiện nguy cơ vượt khả năng kiểm soát. (Vì): Khi phát hiện bất kỳ nguy cơ nào vượt quá khả năng kiểm soát của bản thân, việc báo cáo ngay lập tức cho giáo viên hoặc người phụ trách là yêu cầu bắt buộc để nhận được sự hỗ trợ kịp thời
\item (Đúng) Chỉ tiếp tục thí nghiệm sau khi nguyên nhân nguy hiểm đã được loại bỏ hoặc kiểm soát. (Vì): Thí nghiệm chỉ được phép tiếp tục sau khi nguyên nhân gây ra nguy hiểm đã được loại bỏ hoàn toàn hoặc được kiểm soát chặt chẽ, đảm bảo an toàn cho mọi người.

P: ĐSĐĐ
\end{itemchoice}
}
\end{ex}

% ===== Câu 2 =====
% ID: L10C1B2-01-TL
% Mức: TH
\begin{bt}
% ID: L10C1B2-01-TL
% Mức: TH
Trình bày nguyên tắc cốt lõi khi làm việc với dụng cụ thủy tinh và vật sắc nhọn và giải thích vì sao phải tuân thủ.
\loigiai{
Nguyên tắc cốt lõi là kiểm tra vết nứt, cầm đúng vị trí và thu gom mảnh vỡ bằng dụng cụ chuyên dụng. Phải tuân thủ vì tránh đứt tay, vỡ dụng cụ và gây thương tích thứ cấp; đồng thời luôn dừng và báo người phụ trách khi điều kiện không bảo đảm.
}
\end{bt}

% ===== Câu 3 =====
% ID: L10C1B2-01-TLN
% Mức: NB
\begin{ex}
% ID: L10C1B2-01-TLN
% Mức: NB
Một quy trình về dụng cụ thủy tinh và vật sắc nhọn gồm: (1) nhận diện nguy cơ; (2) kiểm tra vết nứt, cầm đúng vị trí và thu gom mảnh vỡ bằng dụng cụ chuyên dụng; (3) dùng tay không nhặt mảnh kính vỡ; (4) báo giáo viên khi có bất thường; (5) chỉ tiếp tục khi nguy cơ đã được kiểm soát. Có bao nhiêu bước phù hợp quy tắc an toàn? Chỉ ghi số.
\shortans{4}
\loigiai{
Các bước (1), (2), (4), (5) phù hợp; bước (3) không an toàn. Có 4 bước phù hợp.
}
\end{ex}

% ===== Câu 4 =====
% ID: L10C1B2-01-TN
% Mức: NB
\begin{ex}
% ID: L10C1B2-01-TN
% Mức: NB
Trước khi bắt đầu thí nghiệm liên quan đến dụng cụ thủy tinh và vật sắc nhọn, hành động nào an toàn nhất?
\choice
{\True kiểm tra vết nứt, cầm đúng vị trí và thu gom mảnh vỡ bằng dụng cụ chuyên dụng}
{dùng tay không nhặt mảnh kính vỡ}
{tiếp tục thao tác rồi báo cáo sau}
{nhờ bạn khác làm thay mà không kiểm tra điều kiện an toàn}
\loigiai{
Hành động đúng là “kiểm tra vết nứt, cầm đúng vị trí và thu gom mảnh vỡ bằng dụng cụ chuyên dụng” vì tránh đứt tay, vỡ dụng cụ và gây thương tích thứ cấp. Chọn A.
}
\end{ex}

% ===== Câu 5 =====
% ID: L10C1B2-02-DS
% Mức: NB
\begin{ex}
% ID: L10C1B2-02-DS
% Mức: NB
Trong tình huống khi một bạn chưa đọc hướng dẫn đối với dụng cụ thủy tinh và vật sắc nhọn, hãy đánh giá các phát biểu sau.
\choiceTF
{\True Cần kiểm tra vết nứt, cầm đúng vị trí và thu gom mảnh vỡ bằng dụng cụ chuyên dụng.}
{Có thể dùng tay không nhặt mảnh kính vỡ nếu thao tác diễn ra nhanh.}
{\True Phải báo cho giáo viên hoặc người phụ trách khi phát hiện nguy cơ vượt khả năng kiểm soát.}
{\True Chỉ tiếp tục thí nghiệm sau khi nguyên nhân nguy hiểm đã được loại bỏ hoặc kiểm soát.}
\loigiai{
\begin{itemchoice}
\item (Đúng) Cần kiểm tra vết nứt, cầm đúng vị trí và thu gom mảnh vỡ bằng dụng cụ chuyên dụng. (Vì): Cần kiểm tra kỹ dụng cụ thủy tinh xem có vết nứt hay hư hỏng không. Khi sử dụng, phải cầm đúng vị trí để tránh bị trượt tay. Nếu không may làm vỡ, cần thu gom mảnh vỡ cẩn thận bằng các dụng cụ chuyên dụng như kẹp gắp, chổi và hốt rác, không dùng tay không
\item (Sai) Có thể dùng tay không nhặt mảnh kính vỡ nếu thao tác diễn ra nhanh. (Vì): Tuyệt đối không được dùng tay không để nhặt mảnh kính vỡ, dù thao tác có diễn ra nhanh hay chậm. Mảnh kính vỡ rất sắc bén và có thể gây đứt tay nghiêm trọng, dẫn đến chảy máu và nhiễm trùng
\item (Đúng) Phải báo cho giáo viên hoặc người phụ trách khi phát hiện nguy cơ vượt khả năng kiểm soát. (Vì): Nếu phát hiện bất kỳ nguy cơ nào vượt quá khả năng kiểm soát của bản thân, ví dụ như một lượng lớn mảnh vỡ hoặc một tình huống có thể gây nguy hiểm cho người khác, việc đầu tiên cần làm là ngay lập tức báo cáo cho giáo viên hoặc người phụ trách phòng thí nghiệm
\item (Đúng) Chỉ tiếp tục thí nghiệm sau khi nguyên nhân nguy hiểm đã được loại bỏ hoặc kiểm soát. (Vì): Thí nghiệm chỉ được phép tiếp tục sau khi nguyên nhân gây ra nguy hiểm đã được loại bỏ hoàn toàn hoặc được kiểm soát một cách an toàn. Điều này đảm bảo an toàn cho tất cả mọi người tham gia thí nghiệm.

P: ĐSĐĐ
\end{itemchoice}
}
\end{ex}

% ===== Câu 6 =====
% ID: L10C1B2-02-TL
% Mức: TH
\begin{bt}
% ID: L10C1B2-02-TL
% Mức: TH
Phân tích hành vi “dùng tay không nhặt mảnh kính vỡ” và đề xuất cách thay thế an toàn.
\loigiai{
Hành vi “dùng tay không nhặt mảnh kính vỡ” có thể gây tai nạn vì làm mất kiểm soát nguồn nguy hiểm. Cần dừng thao tác, cô lập nguy cơ và thay bằng biện pháp: kiểm tra vết nứt, cầm đúng vị trí và thu gom mảnh vỡ bằng dụng cụ chuyên dụng.
}
\end{bt}

% ===== Câu 7 =====
% ID: L10C1B2-03-TL
% Mức: VD
\begin{bt}
% ID: L10C1B2-03-TL
% Mức: VD
Xây dựng một bảng kiểm ngắn gồm ít nhất bốn mục cho tình huống khi khu vực làm việc bị ướt liên quan đến dụng cụ thủy tinh và vật sắc nhọn.
\loigiai{
Bảng kiểm gồm: đọc hướng dẫn; nhận diện mối nguy; kiểm tra dụng cụ và bảo hộ; thực hiện kiểm tra vết nứt, cầm đúng vị trí và thu gom mảnh vỡ bằng dụng cụ chuyên dụng; theo dõi bất thường; thu dọn và báo cáo sau thí nghiệm.
}
\end{bt}

% ===== Câu 8 =====
% ID: L10C1B2-03-TN
% Mức: NB
\begin{ex}
% ID: L10C1B2-03-TN
% Mức: NB
Trong lúc thay đổi cách bố trí dụng cụ liên quan đến dụng cụ thủy tinh và vật sắc nhọn, hành động nào an toàn nhất?
\choice
{tiếp tục thao tác rồi báo cáo sau}
{nhờ bạn khác làm thay mà không kiểm tra điều kiện an toàn}
{\True kiểm tra vết nứt, cầm đúng vị trí và thu gom mảnh vỡ bằng dụng cụ chuyên dụng}
{dùng tay không nhặt mảnh kính vỡ}
\loigiai{
Hành động đúng là “kiểm tra vết nứt, cầm đúng vị trí và thu gom mảnh vỡ bằng dụng cụ chuyên dụng” vì tránh đứt tay, vỡ dụng cụ và gây thương tích thứ cấp. Chọn C.
}
\end{ex}

% ===== Câu 9 =====
% ID: L10C1B2-04-TL
% Mức: VD
\begin{bt}
% ID: L10C1B2-04-TL
% Mức: VD
Một học sinh đã dùng tay không nhặt mảnh kính vỡ. Hãy nêu trình tự xử lí ngay sau khi phát hiện hành vi này.
\loigiai{
Trình tự: yêu cầu dừng ngay; không tiếp cận nếu chưa an toàn; cô lập nguồn nguy hiểm; báo giáo viên; sơ cứu hoặc xử lí theo hướng dẫn; chỉ hoạt động lại sau khi được xác nhận an toàn.
}
\end{bt}

% ===== Câu 10 =====
% ID: L10C1B2-05-DS
% Mức: VD
\begin{ex}
% ID: L10C1B2-05-DS
% Mức: VD
Trong tình huống khi có người không đeo phương tiện bảo hộ đối với dụng cụ thủy tinh và vật sắc nhọn, hãy đánh giá các phát biểu sau.
\choiceTF
{\True Cần kiểm tra vết nứt, cầm đúng vị trí và thu gom mảnh vỡ bằng dụng cụ chuyên dụng.}
{Có thể dùng tay không nhặt mảnh kính vỡ nếu thao tác diễn ra nhanh.}
{\True Phải báo cho giáo viên hoặc người phụ trách khi phát hiện nguy cơ vượt khả năng kiểm soát.}
{\True Chỉ tiếp tục thí nghiệm sau khi nguyên nhân nguy hiểm đã được loại bỏ hoặc kiểm soát.}
\loigiai{
\begin{itemchoice}
\item (Đúng) Cần kiểm tra vết nứt, cầm đúng vị trí và thu gom mảnh vỡ bằng dụng cụ chuyên dụng. (Vì): Việc kiểm tra vết nứt, cầm đúng vị trí và thu gom mảnh vỡ bằng dụng cụ chuyên dụng là các biện pháp cần thiết để đảm bảo an toàn khi xử lý dụng cụ thủy tinh bị vỡ
\item (Sai) Có thể dùng tay không nhặt mảnh kính vỡ nếu thao tác diễn ra nhanh. (Vì): Tuyệt đối không được dùng tay không để nhặt mảnh kính vỡ, dù thao tác có nhanh đến đâu, vì có thể gây đứt tay nghiêm trọng
\item (Đúng) Phải báo cho giáo viên hoặc người phụ trách khi phát hiện nguy cơ vượt khả năng kiểm soát. (Vì): Khi phát hiện nguy cơ mất an toàn vượt quá khả năng kiểm soát của bản thân, việc báo cáo ngay cho giáo viên hoặc người phụ trách là quy tắc bắt buộc để có biện pháp xử lý kịp thời
\item (Đúng) Chỉ tiếp tục thí nghiệm sau khi nguyên nhân nguy hiểm đã được loại bỏ hoặc kiểm soát. (Vì): Thí nghiệm chỉ được phép tiếp tục khi nguyên nhân gây nguy hiểm đã được loại bỏ hoàn toàn hoặc được kiểm soát chặt chẽ, đảm bảo an toàn cho mọi người.

P: ĐĐĐĐ
\end{itemchoice}
}
\end{ex}

% ===== Câu 11 =====
% ID: L10C1B2-05-TN
% Mức: TH
\begin{ex}
% ID: L10C1B2-05-TN
% Mức: TH
Khi một bạn chưa đọc hướng dẫn liên quan đến dụng cụ thủy tinh và vật sắc nhọn, hành động nào an toàn nhất?
\choice
{\True kiểm tra vết nứt, cầm đúng vị trí và thu gom mảnh vỡ bằng dụng cụ chuyên dụng}
{dùng tay không nhặt mảnh kính vỡ}
{tiếp tục thao tác rồi báo cáo sau}
{nhờ bạn khác làm thay mà không kiểm tra điều kiện an toàn}
\loigiai{
Hành động đúng là “kiểm tra vết nứt, cầm đúng vị trí và thu gom mảnh vỡ bằng dụng cụ chuyên dụng” vì tránh đứt tay, vỡ dụng cụ và gây thương tích thứ cấp. Chọn A.
}
\end{ex}

% ===== Câu 17 =====
% ID: L10C1B2-08-TN
% Mức: VD
\begin{ex}
% ID: L10C1B2-08-TN
% Mức: VD
Khi có người không đeo phương tiện bảo hộ liên quan đến dụng cụ thủy tinh và vật sắc nhọn, hành động nào an toàn nhất?
\choice
{nhờ bạn khác làm thay mà không kiểm tra điều kiện an toàn}
{\True kiểm tra vết nứt, cầm đúng vị trí và thu gom mảnh vỡ bằng dụng cụ chuyên dụng}
{dùng tay không nhặt mảnh kính vỡ}
{tiếp tục thao tác rồi báo cáo sau}
\loigiai{
Hành động đúng là “kiểm tra vết nứt, cầm đúng vị trí và thu gom mảnh vỡ bằng dụng cụ chuyên dụng” vì tránh đứt tay, vỡ dụng cụ và gây thương tích thứ cấp. Chọn B.
}
\end{ex}

% ===== Câu 19 =====
% ID: L10C1B2-09-TN
% Mức: VD
\begin{ex}
% ID: L10C1B2-09-TN
% Mức: VD
Khi kết quả đo khác dự kiến liên quan đến dụng cụ thủy tinh và vật sắc nhọn, hành động nào an toàn nhất?
\choice
{\True kiểm tra vết nứt, cầm đúng vị trí và thu gom mảnh vỡ bằng dụng cụ chuyên dụng}
{dùng tay không nhặt mảnh kính vỡ}
{tiếp tục thao tác rồi báo cáo sau}
{nhờ bạn khác làm thay mà không kiểm tra điều kiện an toàn}
\loigiai{
Hành động đúng là “kiểm tra vết nứt, cầm đúng vị trí và thu gom mảnh vỡ bằng dụng cụ chuyên dụng” vì tránh đứt tay, vỡ dụng cụ và gây thương tích thứ cấp. Chọn A.
}
\end{ex}

% ===== Câu 21 =====
% ID: L10C1B2-10-TN
% Mức: VD
\begin{ex}
% ID: L10C1B2-10-TN
% Mức: VD
Khi giáo viên yêu cầu dừng thao tác liên quan đến dụng cụ thủy tinh và vật sắc nhọn, hành động nào an toàn nhất?
\choice
{dùng tay không nhặt mảnh kính vỡ}
{tiếp tục thao tác rồi báo cáo sau}
{nhờ bạn khác làm thay mà không kiểm tra điều kiện an toàn}
{\True kiểm tra vết nứt, cầm đúng vị trí và thu gom mảnh vỡ bằng dụng cụ chuyên dụng}
\loigiai{
Hành động đúng là “kiểm tra vết nứt, cầm đúng vị trí và thu gom mảnh vỡ bằng dụng cụ chuyên dụng” vì tránh đứt tay, vỡ dụng cụ và gây thương tích thứ cấp. Chọn D.
}
\end{ex}

